% page 101 du fac-simile (image p-100 du scan)
% bloc 162.6 x 233.3 mm ; marge gauche 39.4 mm
\pgc{17.780mm}{0.000mm}{
\l{\cel{52}93}
\l{}
\l{\sou{cirkonspekta}. Qua surveyas prudente, e desfideme, to quon}
\l{\cel{3}lu dicas od agas, o to quon dicas od agas la kunhomi.}
\l{\cel{3}- EFIS}
\l{}
\l{\sou{cirkonstanco}. I. Singla de la fakti di eventajo o di situeso.}
\l{\cel{3}- II. Fakto acesora, qua akompanas eventajo precipua. - EFIS.}
\l{}
\l{\sou{cirkonvalaciono}. (\sou{milit.}) Trancheo kun palisadi, parapeti, for-}
\l{\cel{3}tifikajo tempala per qua armeo envelopas la placo +asiejata,}
\l{\cel{3}por cesigar la relati di olca kun la mondo extera, o per qua}
\l{\cel{3}lu su protektas kontre la ataki de la trupi qui venas por}
\l{\cel{3}sokursar la placo +asiejata. - EFIRS.}
\l{}
\l{\sou{cirkonvoluciono}.\cel{2}Envolveso - sinui cirkloforma. -\cel{2}EFIS.}
\l{}
\l{\sou{cirkuito}. I. Spaco parkurenda por cirkumirar komplete loko.}
\l{\cel{3}- II. Iro cirkuma, vice iro tra la voyo rekta. - EFIS.}
\l{}
\l{\sou{cirkular}. (\sou{netrans}.) I. Irar e retroirar omna-direcione sinue.}
\l{\cel{3}II. (\sou{pri moneto}).\cel{2}Aceptesar da omna homi di la regiono a}
\l{\cel{3}qua lu esas destinita, kom valoroza. - DEFIRS.}
\l{}
\l{\sou{cirkulero}. Letro de qua un kopiuro uniforma sendesas sama-}
\l{\cel{3}instante a plura (e, generale : a tre multa) personi. -}
\l{\cel{3}DEFIRS.}
\l{}
\l{\sou{cirkum}. En la loko, tempo o quanto qua esas vicina, apuda}
\l{\cel{3}- kompare a la objekto, quaze ringe exter lu. - DL.}
\l{}
\l{\sou{cirkumnutaciono}. (\sou{bot}.) Parkuro helicoida da la extremajo di}
\l{\cel{3}la radiko, o di la stipo, di vejetanto.}
\l{}
\l{\sou{\textquotedbl{}text}quotedbl\{\}cis\textquotedbl{} . (Muziko).\textquotedbl{}Do\textquotedbl{}-noto, diezoza.}
\l{}
\l{\sou{cis.}\cel{3}Ca-latere di - ne trans lu. - L.}
\l{}
\l{\sou{cisipara}. (\sou{biol.}) Qua su riproduktas per sua divideso a du}
\l{\cel{3}parti, generale segun la mikra axo di la korpo. - DEFIRS.}
\l{}
\l{\sou{cisoido}. (\sou{geom.}) Kurvo di la grado triesma, konsistanta ye du}
\l{\cel{3}branchi simetra qui departas de un punto e tendencas a tan-}
\l{\cel{3}gento komuna, sen atingar lu ul-instante. - DEF.}
\l{}
\l{\sou{cisterciano}. (\sou{religio kristana)}.\cel{2}Reguliero qua resortias a}
\l{\cel{3}la ordeno religiala qua kreesis en la Franca urbo Citeaux.}
\l{\cel{3}DEFIRS.}
\l{}
\l{\sou{cisterno}. Granda tanko a qua duktesas e rekoliesas la aquo}
\l{\cel{3}pluvala. - DEFIRS.}
\l{}
\l{\sou{citar}. (\sou{trans.}) Adportar voce o skribe extraktaji de texto}
\l{\cel{3}por sekondar to quon onu asertas. - DEFIRS.}
}
